\documentclass[11pt,a4j,twocolumn]{jarticle}
\usepackage{graphicx}

% 設問(a)に関する解答はこれより下に記すこと。
\title{情報リテラシー}
\author{1821013　岩谷　慶}
\date{\today}
% 設問(a)に関する解答はこれより上に記すこと。

\begin{document}
\maketitle
\section{箇条書きの問題}

% 設問(b)に関する解答はこれより下に記すこと。
知識工学部には以下のような学科があります。
\begin{itemize}
\item 情報科学
\item 情報通信学
\item 経営システム工学
\item 自然科学
\end{itemize}
% 設問(b)に関する解答はこれより上に記すこと。

\section{数式の問題}
% 設問(c)に関する解答はこれより下に記すこと。
2次方程式 $ ax^{2}+bx+c=0 $の解は次式で与えられます。
\begin{eqnarray}
  x_{1} & = & \frac{-b + \sqrt{b^2 - 4ac}}{2a} \\
  x_{2} & = & \frac{-b + \sqrt{b^2 - 4ac}}{2a}
  %\sqrt{y+z}
\end{eqnarray}

% 設問(c)に関する解答はこれより上に記すこと。

\section{図の貼付の問題}

% 設問(d)に関する解答はこれより下に記すこと。
図\ref{figure1}に示すのは、東京都市大学のマスコットです。
\begin{figure}[t]
\begin{center}
\includegraphics[width=55mm]{mascot.eps}
\end{center}
\caption{大学マスコット}
\label{figure1}
\end{figure}

% 設問(d)に関する解答はこれより上に記すこと。

\section{表作成の問題}
% 設問(e)に関する解答はこれより下に記すこと。
表\ref{table1}に示すのは、卒業までの単位取得目安数です。

\begin{table}[h]
\caption{単位取得目安}
\label{table1}
\begin{center}
\begin{tabular}{|l|c|c|c|c|}\hline
     &１年　&　２年　&　３年　&　４年 \\ \hline
    単位数 & ４０ & ８０ & １２０ & １２６ \\  \hline 
\end{tabular}
\end{center}
\end{table}

% 設問(e)に関する解答はこれより上に記すこと。

\section{参考文献の問題}
% 設問(f)に関する解答はこれより下に記すこと。
東京都市大学にはYouTubeに公式アカウント\cite{ref1}があります。

\begin{thebibliography}{9}
\bibitem{ref1}都市大公式,https://www.youtube.com/tcu104
\end{thebibliography}

% 設問(f)に関する解答はこれより上に記すこと。

\end{document}
